\documentclass{butexFR}


\titre{Solution aux besoins}
\UE{Analyse}
\sujet{SAÉ 3.02 - Labyrinthe}
\enseignant{Mme Boneva}
\eleves{Bredelle Victor \\ Marouze Antonin \\ Harlaut Romain \\ Zheng Angèl \\ Lavogiez Baptiste}


\begin{document}

\fairepagedegarde
\fairemarges
\tabledematieres



\section{Introduction}

Le présent rapport détaille la composante \texttt{Analyse} relative au projet \textbf{SAÉ S3.02 - Labyrinthe}. Il aborde l'organisation interne de notre équipe ainsi que notre approche de modélisation des exigences du client, structurée comme suit :

\begin{enumerate}
    \item Diagrammes de cas d'utilisation
    \item Fiches descriptives
    \item Prototype d'interface
    \item Diagramme de classe
\end{enumerate}

Il est entendu que chaque section est itérative.

\section{Équipe et organisation du travail}
\textbf{Groupe G3.3}

\begin{itemize}
    \item Bredelle Victor
    \item Marouze Antonin
    \item Harlaut Romain
    \item Zheng Angèl
    \item Lavogiez Baptiste
\end{itemize}

\tsec{Victor}

Rédaction de trois fiches descriptives avec Baptiste :
\begin{enumerate}
    \item Charger profil
    \item Afficher progression
    \item Lancer mode progression
\end{enumerate}

\tsec{Antonin}

Réalisation de la maquette \textit{Figma} avec % ⟨1: Original: "Angèl / Romain"⟩
Angèl et Romain.

\tsec{Romain}

Réalisation de la fiche descriptive "Avancer d'une case", aide à la maquette \textit{Figma}.

\tsec{Angèl}

Réalisation de la maquette \textit{Figma} avec % ⟨2: Original: "Antonin / Romain"⟩
Antonin et Romain.

\tsec{Baptiste}

Rédaction de trois fiches descriptives avec Victor :
\begin{enumerate}
    \item Charger profil
    \item Afficher progression
    \item Lancer mode progression
\end{enumerate}

Rédaction du présent rapport.

\tsecnonum{Le groupe}

Outre les responsabilités individuelles, nous avons collectivement débattu de la manière de modéliser chaque exigence en l'examinant sous divers angles.

\info{\textbf{L'ensemble des contributeurs a collaboré à l'élaboration du diagramme de classes.}}

\section{Cas d'utilisation}

\subsection{Diagrammes de cas d'utilisation}

\begin{figure}[h!]
    \centering
    \includegraphics[width=0.6\textwidth]{figures/cas_dutilisation/diagrammes_de_cas_dutilisation/default/fig_1.png}
    \caption{Cas d'utilisation identifiés}
    \label{fig:cas_dutilisation_diagrammes_de_cas_dutilisation_1}
\end{figure} /

(Review en TP avec autres)

\newpage

\subsection{Fiches descriptives}

\subsubsection{% ⟨3: Original: "Cas d’utilisation: Charger profil"⟩
Cas d’utilisation : Charger profil}
\textbf{Système:} Labyrinthe\\
\textbf{Acteur principal:} Utilisateur\\
\textbf{Déclencheur:} /\\
\textbf{Autres acteurs:} /\\[0.25em]
\textbf{Préconditions:} L'utilisateur est dans le menu de saisie du nom (``Lancer le mode progression'')\\
\textbf{Garanties en cas de succès:} Le système charge le profil associé au nom et l'utilisateur reprend sa progression\\
\textbf{Garanties minimales:} /\\[0.25em]
\textbf{Scénario nominal:}
\begin{enumerate}
    \item L'utilisateur a rentré un nom déjà existant dans la base de données lors du \textbf{Cas d'utilisation} ``Lancer le mode progression'' (suite de ce cas, ayant déjà vérifié % ⟨7: Original: "l'occurence"⟩
l'occurrence)
    \item Le système affiche un menu où l'utilisateur peut choisir de ``Reprendre la partie existante'' ou ``Écraser la partie existante''.
    \item L'utilisateur sélectionne ``Reprendre la partie''
    \item Le système charge le profil associé au nom depuis la base de données et l'utilisateur reprend sa progression
\end{enumerate}
\textbf{Scénario alternatifs:} /

\newpage

\subsubsection{% ⟨9: Original: "Cas d’utilisation: Afficher progression"⟩
Cas d’utilisation : Afficher progression}
\textbf{Système:} Labyrinthe\\
\textbf{Acteur principal:} Utilisateur\\
\textbf{Déclencheur:} /\\
\textbf{Autres acteurs:} /\\[0.25em]
\textbf{Préconditions:} % ⟨12: Original: "Avoir un profil, être dans le mode progression"⟩
L'utilisateur possède un profil et se trouve en mode progression.\\
\textbf{Garanties en cas de succès:} Le système affiche la progression associée au joueur (score, caractéristiques...)\\
\textbf{Garanties minimales:} /\\[0.25em]
\textbf{Scénario nominal:}
\begin{enumerate}
    \item L'utilisateur sélectionne la fonctionnalité ``accéder au menu du mode progression''
    \item Le système récupère les données du joueur dans la base et % ⟨14: Original: "calcule le score et caractéristiques"⟩
calcule le score et les caractéristiques du joueur afin d'afficher la progression
\end{enumerate}
\textbf{Scénario alternatifs:} /

\newpage

\subsubsection{% ⟨16: Original: "Cas d’utilisation: Lancer mode progression"⟩
Cas d’utilisation : Lancer mode progression}
\textbf{Système:} Labyrinthe\\
\textbf{Acteur principal:} Utilisateur\\
\textbf{Déclencheur:} /\\
\textbf{Autres acteurs:} /\\[0.25em]
\textbf{Préconditions:} % ⟨19: Original: "/ (avoir lancé le jeu)"⟩
Le jeu a été lancé.\\
\textbf{Garanties en cas de succès:} L'utilisateur % ⟨20: Original: "arrive sur le mode progression"⟩
accède au mode progression\\
\textbf{Garanties minimales:} /\\[0.25em]
\textbf{Scénario nominal:}
\begin{enumerate}
    \item L'utilisateur sélectionne la fonctionnalité ``Lancer mode progression''.
    \item Le système affiche un menu où l'utilisateur peut % ⟨22: Original: "rentrer son nom"⟩
saisir son nom
    \item L'utilisateur % ⟨23: Original: "rentre son nom"⟩
saisit son nom
    \item Le système vérifie que le nom est valide, sauvegarde son profil dans la base de données et lance le mode progression
\end{enumerate}
\textbf{Scénario alternatifs:}

\emph{A: Le système détecte une erreur lors de la vérification}\\
\hspace*{1em}4(A): Le système affiche une notification % ⟨24: Original: "(de quelles mauvaises informations il s'agit)"⟩
(précisant la nature de l'erreur), puis revient en étape 2 du scénario principal

\emph{B: Le système détecte une occurence du nom dans la base de données (charger partie)}\\
\hspace*{1em}4(B): Le système démarre le \textbf{Cas d'utilisation} ``Charger profil''

\newpage

\subsubsection{% ⟨26: Original: "Cas d’utilisation: Déplacer joueur"⟩
Cas d’utilisation : Déplacer joueur}
\textbf{Système:} Labyrinthe\\
\textbf{Acteur principal:} Utilisateur\\
\textbf{Déclencheur:} /\\
\textbf{Autres acteurs:} /\\[0.25em]
\textbf{Préconditions:} % ⟨29: Original: "Être dans le labyrinthe"⟩
Le joueur se trouve dans le labyrinthe.\\
\textbf{Garanties en cas de succès:} Le système % ⟨30: Original: "a déplacé le joueur"⟩
déplace le joueur à la case souhaitée\\
\textbf{Garanties minimales:} /\\[0.25em]
\textbf{Scénario nominal:}
\begin{enumerate}
    \item Le joueur appuie sur une flèche directionnelle
    \item Le système vérifie que la direction est accessible et déplace le joueur sur la case associée
\end{enumerate}
\textbf{Scénario alternatifs:}

\emph{A: Le joueur % ⟨32: Original: "arrive sur la sortie"⟩
atteint la sortie}\\
\hspace*{1em}2(A): Le système affiche une notification annonçant que la sortie a été atteinte et redirige vers le menu principal.

\emph{B: Le joueur veut se déplacer dans une case inaccessible}\\
\hspace*{1em}2(B): Le système affiche la position actuelle % ⟨33: Original: "et pas la nouvelle position"⟩
et non la nouvelle position et revient en étape 1 du scénario principal

\newpage

\subsection{Prototypes d'interface}

\begin{figure}[h!]
    \centering
    \includegraphics[width=1\textwidth]{figures/cas_dutilisation/prototypes_dinterface/default/fig_1.png}
    \caption{Prototype principal}
    \label{fig:cas_dutilisation_prototypes_dinterface_1}
\end{figure}

% ⟨34: Original: "Faire photos successives ?"⟩
Il serait pertinent d'ajouter des photos successives pour illustrer les étapes.

L'interface est pensée pour être simple, intuitive, et sans % ⟨35: Original: "ambiguité"⟩
ambiguïté. La priorité demeure, comme toujours, le \textbf{strict respect des attentes du client}.

Ces prototypes ont été élaborés préalablement à la rédaction des fiches descriptives, dans le but de garantir une modélisation adéquate de la séquence, conformément à la méthodologie enseignée dans le cadre du cours d'\texttt{Analyse}.

\section{Diagramme de classe}

\begin{figure}[h!]
    \centering
    \includegraphics[width=1\textwidth]{figures/diagramme_de_classe/default/default/fig_1.png}
    \caption{Diagramme de classe identifié}
    \label{fig:diagramme_de_classe_default_1}
\end{figure}

Le diagramme de classe a été généré à l'aide de \textit{PlantUML} dans un format textuel, assurant ainsi une portabilité et une facilité de modification accrues, quelle que soit la localisation.

La conception de ce diagramme a nécessité un investissement conséquent, considérant que, par nature, un diagramme de classe requiert une \textbf{réflexion soutenue et l'intégration de multiples perspectives}. 

Des retours ont été sollicités auprès des professeurs spécialisés en analyse, ainsi qu'auprès de ceux enseignant la qualité de développement.

\newpage

Par conséquent, nous estimons avoir abouti à une solution hautement optimisée, caractérisée par l'absence de redondance et le respect du \textbf{principe de responsabilité unique} pour chaque classe. 

\merci

\end{document}




